% !TeX root = ../../../navier-stokes-manuscript.tex
% ============================================================
% 2.3 The Pressure-Complete Mixed Jet
% ============================================================

\subsection{The Pressure-Complete Mixed Jet}\label{sub:ThePressureCompleteMixedJet}

Return to the original incompressible system
\[
  \partial_tu+(u\cdot\nabla)u+\nabla p
  =\nu\Delta u,
  \qquad
  \nabla\cdot u=0.
\]
Every derivative relation in this section will come from this equation itself.
For a temporal order \(m\in\Nzero\) and a spatial multi-index \(\alpha\in\Nzero^3\), define
\[
  J_{m,\alpha}
  :=
  \partial_t^m\partial_x^\alpha u,
  \qquad
  \Pi_{m,\alpha}
  :=
  \partial_t^m\partial_x^\alpha p.
\]
The object \(J_{m,\alpha}\) is one velocity coordinate of the mixed jet, and \(\Pi_{m,\alpha}\) is the pressure coordinate attached to it.
The pair of indices preserves the independent temporal and spatial depths exposed by the finite-time response argument.

The low rows show how the nesting begins.
Differentiate the equation once in time.
Leibniz's rule splits the derivative of transport across the velocity that carries and the velocity that is carried:
\[
  \partial_t\bigl((u\cdot\nabla)u\bigr)
  =
  (\partial_tu\cdot\nabla)u
  +(u\cdot\nabla)\partial_tu.
\]
The first differentiated temporal row is
\[
  J_{2,0}
  +(J_{1,0}\cdot\nabla)J_{0,0}
  +(J_{0,0}\cdot\nabla)J_{1,0}
  +\nabla\Pi_{1,0}
  =
  \nu\Delta J_{1,0}.
\]
The next time response \(J_{2,0}\) appears beside two interactions between the velocity and its first time response.
Pressure differentiates at the same temporal depth, and viscosity acts on the spatial curvature of the row being evolved.

Now differentiate once in a spatial direction \(x_k\).
The transport term splits in the same way:
\[
  \partial_k\bigl((u\cdot\nabla)u\bigr)
  =
  (\partial_ku\cdot\nabla)u
  +(u\cdot\nabla)\partial_ku.
\]
Writing \(e_k\) for the corresponding spatial multi-index gives
\[
  J_{1,e_k}
  +(J_{0,e_k}\cdot\nabla)J_{0,0}
  +(J_{0,0}\cdot\nabla)J_{0,e_k}
  +\nabla\Pi_{0,e_k}
  =
  \nu\Delta J_{0,e_k}.
\]
The spatial derivative is transported by the existing velocity and also changes the velocity doing the transporting.
The same row retains pressure and the two additional spatial derivatives carried by the Laplacian.

The general relation follows by applying \(\partial_t^m\partial_x^\alpha\) to the same equation.
For multi-indices \(\beta\leq\alpha\), write
\[
  \binom{\alpha}{\beta}
  :=
  \prod_{j=1}^{3}\binom{\alpha_j}{\beta_j}.
\]
The double Leibniz rule gives
\[
  \partial_t^m\partial_x^\alpha\bigl((u\cdot\nabla)u\bigr)
  =
  \sum_{a=0}^{m}
  \sum_{\beta\leq\alpha}
  \binom{m}{a}
  \binom{\alpha}{\beta}
  (J_{a,\beta}\cdot\nabla)
  J_{m-a,\alpha-\beta}.
\]
Substitution gives the full pressure-complete recurrence
\begin{align*}
  J_{m+1,\alpha}
  &+
  \sum_{a=0}^{m}
  \sum_{\beta\leq\alpha}
  \binom{m}{a}
  \binom{\alpha}{\beta}
  (J_{a,\beta}\cdot\nabla)
  J_{m-a,\alpha-\beta}
  \\
  &+
  \nabla\Pi_{m,\alpha}
  =
  \nu\Delta J_{m,\alpha},
  \qquad
  \nabla\cdot J_{m,\alpha}=0.
\end{align*}
Every way of distributing the temporal and spatial derivatives across the two velocity factors appears once with its binomial multiplicity.
The recurrence is triangular in temporal order because the new coordinate \(J_{m+1,\alpha}\) is expressed through rows at temporal depth at most \(m\), together with their spatial derivatives and pressure response.
It is simultaneously open in spatial height because the transport gradient and viscous Laplacian reach beyond \(\alpha\).

Pressure remains part of the construction at every row.
Taking the divergence of the original momentum equation and using incompressibility gives
\[
  -\Delta p
  =
  \partial_i\partial_j(u_i u_j).
\]
Applying the same mixed derivative yields
\[
  -\Delta\Pi_{m,\alpha}
  =
  \partial_i\partial_j
  \partial_t^m\partial_x^\alpha(u_i u_j).
\]
Expanding the product makes its dependence on the velocity jet explicit:
\[
  -\Delta\Pi_{m,\alpha}
  =
  \partial_i\partial_j
  \sum_{a=0}^{m}
  \sum_{\beta\leq\alpha}
  \binom{m}{a}
  \binom{\alpha}{\beta}
  (J_{a,\beta})_i
  (J_{m-a,\alpha-\beta})_j.
\]
Solving this Poisson equation is nonlocal, so the pressure at one point depends on the velocity products across the whole spatial field.
Pressure cannot be replaced by an independent local neighbour edge or discarded as an external forcing.
It is the incompressible response generated by the same mixed velocity coordinates.

The recurrence reveals the complete derivative object behind the response tower.
Time differentiation, spatial differentiation, transport, viscous spatial shift, incompressibility, and nonlocal pressure all remain joined.
It also reveals a new difficulty: the full tensor relation is too large to show plainly how each temporal advance opens a higher spatial Sobolev row.
The critical height supplies a scalar window in which that nesting can be derived without losing the lower production terms that made it.
